% ============================================================
% f2014_anhui_science_q13_point_plot.tex
% 2014 安徽卷（理科）第 13 题：点 A_i(i,a_i) 的坐标图
% 来源 PDF：普通高考/2014/2014安徽理.pdf
% 原页：第 1 页（A3 横排）中栏下部，印面 857，第 13 题题干右侧
% 说明：A_0=(0,1)，A_1=(1,3)，A_2=(2,4)；虚线为向两坐标轴的投影。
% ============================================================
\begin{tikzpicture}[>=Stealth, line width=0.7pt, font=\small,
  x=1.15cm, y=1.15cm, line cap=round, line join=round]

  % 坐标轴
  \draw[->] (-0.45,0) -- (3.35,0) node[right] {$x$};
  \draw[->] (0,-0.45) -- (0,4.85) node[above] {$y$};
  \node[below left] at (0,0) {$O$};

  % 刻度
  \foreach \x in {1,2} {
    \draw (\x,0.06) -- (\x,-0.06);
    \node[below] at (\x,-0.08) {$\x$};
  }
  \foreach \y/\lab in {1/1,3/3,4/4} {
    \draw (0.06,\y) -- (-0.06,\y);
    \node[left] at (-0.08,\y) {$\lab$};
  }

  % 投影虚线
  \draw[dashed] (1,3) -- (1,0);
  \draw[dashed] (2,4) -- (2,0);
  \draw[dashed] (1,3) -- (0,3);
  \draw[dashed] (2,4) -- (0,4);

  % 点
  \fill (0,1) circle (1.6pt) node[right, inner sep=2pt] {$A_0$};
  \fill (1,3) circle (1.6pt) node[right, inner sep=2pt] {$A_1$};
  \fill (2,4) circle (1.6pt) node[right, inner sep=2pt] {$A_2$};

\end{tikzpicture}
